\documentclass[pdftex,12pt,a4paper]{report}

\usepackage[portuguese,english]{babel}
\usepackage[T1]{fontenc} 
\usepackage[utf8]{inputenc}
\usepackage[pdftex]{graphicx}
\usepackage{minitoc}
\usepackage{hyperref}
\usepackage{indentfirst}
\usepackage[compact]{titlesec}
\usepackage{fancyhdr}
\usepackage{caption}
\usepackage{pgfplots}
\usepackage{pgfplotstable}
\usepackage{fixltx2e}
\usepackage{mathtools}
\usepackage{fancyhdr}

\pagestyle{fancy}
\renewcommand*\thesection{\thechapter\arabic{section}}
\newcommand{\HRule}{\rule{\linewidth}{0.5mm}}
\begin{document}

\begin{titlepage}

\begin{center}

\includegraphics[width=0.15\textwidth]{./logo}\\[0.5cm]    

\textsc{\large Universidade de Aveiro \\[1cm]\large departamento de electrónica, telecomunicações e informática}\\[1cm]

\textsc{\large{42532}\large - Base de Dados \\[1cm]}

\HRule \\[0.5cm]
{ \huge \bfseries Football Club}\\[0.4cm]
{ \large \bfseries Trabalho Prático Final}\\[0.4cm]
\HRule \\[1cm]

\textsc{\small{8240 - MESTRADO INTEGRADO EM ENGENHARIA DE COMPUTADORES E TELEMÁTICA}}\\[1cm]

\begin{minipage}{0.4\textwidth}

\begin{flushleft} \large
\href{mailto:rafael.ferreira@ua.pt}{António Rafael da \\ Costa Ferreira }
 \small{\\NMec: 67405 | P4G5}
\end{flushleft}
\end{minipage}
\begin{minipage}{0.4\textwidth}

\begin{flushright} \large
\href{mailto:rodrigocunha@ua.pt}{Rodrigo Lopes \\ da Cunha}
\small{\\NMec: 67800 | P4G5}
\end{flushright}
\end{minipage}\\[1cm]

{\large Docente: Carlos Manuel Azevedo Costa   }\\[0.5cm]

\vfill

{\large Junho de 2015 \\ 2014-2015}

\end{center}

\end{titlepage} %Titulo do Relatorio
\renewcommand{\headrulewidth}{0pt}

%Cabeçalhos de rodapé
\fancyhead{}
\fancyfoot{}
\lhead{Football Club - Trabalho Prático Final}
\rhead{BD - 2014/2015}
\lfoot{\textit{P4G5} \\ Rafael Ferreira nmec: 67405 \\ Rodrigo Cunha nmec: 67800}
\rfoot{\thepage}

%Renomear Comandos
\renewcommand*\contentsname{Conteúdos}
\renewcommand*\figurename{Figura}
\renewcommand*\tablename{Tabela}

%Conteúdos, dar paragrafo
\tableofcontents
%Headers
\renewcommand{\headrulewidth}{0.15pt}
\renewcommand{\thechapter}{}

\clearpage

\section{Introdução}
% o que, porquê e o objetivo
O trabalho proposto para o projeto da unidade curricular de Base de Dados é uma plataforma de gestão de um clube de futebol. Usando os conhecimentos adquiridos, propôs-se o desenvolvimento deste projeto visto que o futebol é uma modalidade mundial, envolvendo vários tipos de interesse. 

O objetivo desta base de dados desenvolvida é permitir a gestão de todos os processos de um clube de futebol, como será visto mais à frente.

Além de se ter desenvolvido esta base de dados, desenvolveu-se uma aplicação WPF C\# para permitir a manipulação dos dados da base de dados de forma mais simplificada para um utilizador final.

Esta base de dados deve fornecer ferramentas que permitam a criação, remoção, alteração e consulta da base de dados de forma segura, eficiente e robusta.

O relatório reflete todos os passos e decisões tomadas na criação da base de dados que sustenta o projeto bem como uma descrição das capacidades da aplicação desenvolvida para o cliente.

Para a criação deste projeto foi seguido o processo leccionado nas aulas, sendo estas as seguintes fases do processo: análise de requisitos, desenho conceptual, desenho do esquema lógico, desenho do esquema físico e administração.

\newpage
\section{Análise de Requisitos}

A análise de requisitos foi uma das partes mais importantes do processo de concepção do projeto uma vez que ajudou a ter uma visão clara do que o sistema teria de suportar. Após realizado um "brainstorming", estas são as características que o sistema deve suportar:

Uma pessoa que pode ser sócio ou assalariado, pertence ao clube. Uma pessoa é caracterizada por Nome, Data de nascimento, Nacionalidade, B.I., Endereço e NIF.

Um assalariado tem salário e um ID Interno que o identifica no clube. 

Um sócio tem nº sócio, quotas em dia ou não, valor das quotas do sócio que podem depende de sócio para sócio. Um sócio pode ter um lugar anual para assistir aos jogos. Este lugar anual tem um valor, data início, duração e ID Lugar e este lugar anual senta-se num lugar que é caracterizado por um nº de lugar que pertence a uma Secção e esta secção tem um ID da secção e nº de lugares.
Um assalariado faz parte dos trabalhadores do clube e é uma pessoa. Um assalariado pode ser um Jogador, Técnicos ou Staff do clube. Um jogador tem ID da federação, peso e altura e pertence a um escalão. Um técnico, que é um assalariado, pode ter um ID da federação e função. Staff da clube, que é um assalariado que tem um cargo e que trabalha num Departamento que é caracterizado por um Endereço e ID do Departamento. 



\end{document}
